\uuid{LGC1}
\titre{Produit}
\auteur{bodin ; hanani ; mzari}
\organisation{exo7}
\licence{CC BY-NC-SA 4.0}
\datecreate{2020-01-20}
\niveau{L1}
\module{Algèbre}
\chapitre{Polynôme, fraction rationnelle}
\sousChapitre{Définition, degré, produit}
\difficulte{1}
\theme{Polynômes, produit}
\exo7qcmid{lille-1.polynomes.001}

\contenu{
\begin{question}
\qclassification{105.05}
\qtags{motcle=produit}
Soit $P(X) = 2X^5+3X^2+X$ et $Q(X) = 3X^2-2X+3$.
Quelles sont les assertions vraies concernant le polynôme produit $P(X)\times Q(X)$ ?
\begin{answers}
    \bad{Le coefficient dominant est $5$.}

    \good{Le coefficient du monôme $X^3$ est $-3$.}

    \bad{Le coefficient du terme constant est $3$.}

    \good{Le produit est la somme de $7$ monômes ayant un coefficient non nuls.}
\end{answers}

\begin{explanations}
$P(X)\times Q(X) = 6 X^7 - 4 X^6 + 6 X^5 + 9 X^4 - 3 X^3 + 7 X^2 + 3 X$.
\end{explanations}
\end{question}
}
